\usepackage[T1]{fontenc}
\usepackage[utf8]{inputenc}
\usepackage[ngerman]{babel}
\usepackage{tikz}

% TUC-Templates laden.
\usetheme{tuc2019}
\mode<article>{\usepackage{beamerarticletuc2019}}


% customization
\bibliographystyle{alpha}

% deactivate navigation
\setbeamertemplate{navigation symbols}{}

% metadata
\newcommand{\Title}{Benutzerfreundlichkeit fünf verschiedener Zwei-Faktor-Authentifizierungsmethoden}
\newcommand{\ShortTitle}{Zwei-Faktor-Authentifizierungsmethoden}
\title[\ShortTitle]{\Title}
\subtitle{Hauptseminar Medieninformatik}
\author{Joshua Jeschek}
\date{11.06.2022}
\institute[TUC]{TU Chemnitz}
\titlegraphic{\includegraphics[height=0.19\textheight]{tuc2019/logo/tuc_green}}
\tucurl{https://hauptseminar.jeschek.eu}


\begin{document}
\tucthreeheadlines{}
\frame{\titlepage}

\tuctwoheadlines{}
\section*{\ShortTitle}
\subsection{Einleitung}
\section*{\ShortTitle}
\subsection*{Einleitung}
\begin{frame}
    \begin{center}
        \begin{figure}
            \begin{tikzpicture}
                \begin{scope} [fill opacity=.4]
                    \node[circle,draw,fill=gray,minimum width=4cm,anchor=-35]
                    (know) at (0,1) {
                        \includegraphics[width=0.2\textwidth]{password}};
                    \node[circle,draw,fill=gray,minimum width=4cm,anchor=215]
                    (have) at (0,1) {
                        \includegraphics[height=0.2\textwidth]{phone}};
                    \node[circle,draw,fill=gray,minimum width=4cm,anchor=90]
                    (are) at (0,1) {
                        \includegraphics[height=0.2\textwidth]{touch-id}};
                \end{scope}
            \end{tikzpicture}
            \caption{Verschiedene Möglichkeiten für 2FA
                \endnote{\url{https://www.pngegg.com/en/png-eyxan}}
                \endnote{\url{https://images.vexels.com/media/users/3/157570/isolated/lists/4b39b362c76ea5a00de62f8ff839b5ed-einfaches-smartphone-symbol.png}}
                \endnote{\url{https://cdn4.iconfinder.com/data/icons/apple-touch-id/512/Touch_ID-512.png}}
            }
        \end{figure}
    \end{center}

    \note{
        Intro\ldots
    }
\end{frame}


\subsection{Gliederung}
\frame{\frametitle{Gliederung}\tableofcontents[hideallsubsections]}

%
% CONTENT
%

\section{Fünf Zwei-Faktor-Authentifizierungsmethoden}
\frame{\tableofcontents[sectionstyle=show/hide,subsectionstyle=show/show/hide]}

\subsection{SMS}
\begin{frame}
    \frametitle{SMS}
    \begin{enumerate}
        \item Bestätigungscode wird als SMS-Nachricht gesendet
        \item Eingabe des Codes in App
    \end{enumerate}
    \begin{itemize}
        \item am weitesten verbreitet\cite{reese2019}
        \item einige Usability- und Sicherheitsprobleme, aber einfaches Setup
    \end{itemize}
    \note{
        Usability-Probleme:\cite{reese2019}
        \begin{itemize}
            \item verzögerte Zustellung
            \item kein Empfang
            \item Fehler bei Eingabe des Codes
        \end{itemize}
        Sicherheits-Probleme:\cite{reese2019}
        \begin{itemize}
            \item keine Verschlüsselung -> Man In The Middle
            \item SIM-Swapping
            \item Server muss Code sicher speichern
            \item Brute-Force Angriffe sind nicht ausgeschlossen
        \end{itemize}
    }
\end{frame}

\subsection{TOTP}
\begin{frame}
    \frametitle{TOTP}
    \begin{enumerate}
        \item Nutzer scannt QR-Code, Scanner generiert Bestätigungscode
        \item Bestätigungscode wird eingegeben und vom Server überprüft
    \end{enumerate}
    \begin{itemize}
        \item Code hat beschränkte Lebensdauer
        \item Spezieller Scanner oder Smartphone benötigt
    \end{itemize}
    \note{
        \begin{itemize}
            \item benötigt kein Internet oder Netz
        \end{itemize}
        Usability:\cite{reese2019}
        \begin{itemize}
            \item relative kompliziertes Setup
            \item nicht alle haben Smartphone
            \item Fehler bei Eingabe des Codes
        \end{itemize}
        Sicherheit:\cite{reese2019}
        \begin{itemize}
            \item Verschlüsselung \textrightarrow{} besser als SMS
            \item Server und Client müssen Code sicher speichern
        \end{itemize}
    }
\end{frame}

\subsection{Vor-generierte Codes}

\subsection{Push Benachrichtigungen}

\subsection{U2F Sicherheitsschlüssel}


\tucthreeheadlines{}
\section{}
\frame<beamer| trans>{\begin{center}\scalebox{2}{Vielen Dank für eure Aufmerksamkeit.}\end{center}}

% \begin{onlyenv}<beamer| trans>
% \appendix
% \newcounter{finalframe}
% \setcounter{finalframe}{\value{framenumber}}

% %
% % "Spare-Slides" hier einsetzen.
% %

% \setcounter{framenumber}{\value{finalframe}}
% \end{onlyenv}

\tuctwoheadlines{}
\section*{\ShortTitle}
\subsection{Literatur}
% \begin{frame}
\frametitle{Literatur}

\begin{thebibliography}{xxx}
\bibitem{reese2019}
    Reese, K., Smith, T., Dutson, J., Armknecht, J., Cameron, J. \& Seamons, K., 2019 \---
    A Usability Study of Five Two-Factor Authentication Methods. SOUPS @ USENIX Security Symposium.
\newblock{\url{https://www.usenix.org/conference/soups2019/presentation/reese}}
% \newblock{}
\end{thebibliography}
\end{frame}

\begin{frame}
\frametitle{Literatur}
\bibliography{refs}
\end{frame}

\end{document}
